\subsection{傅里叶余变换}

在集合 \(\widehat{\mathbf{G}}\) 上赋予离散拓扑。记 \(\mathrm{F}(\widehat{\mathbf{G}})\) 为各 \(\operatorname{End}(\mathrm{E}_{\pi})\) 的乘积代数，其中 \(\pi\) 遍历 \(\widehat{\mathbf{G}}\)，记 \(\mathrm{F}_b(\widehat{\mathbf{G}})\) 为各 \(\operatorname{End}(\mathrm{E}_{\pi})\) 的乘积星代数（I, 第 103 页例 5）；这就是由族 \((u_{\pi})_{\pi \in \widehat{\mathbf{G}}}\) 组成且满足 \(\sup \| u_{\pi} \| < +\infty\) 的集合。

\[
\pi \in \widehat {\mathbf {G}}
\]

记 \(\mathrm{F}_{0}(\widehat{\mathrm{G}})\) 为 \(\mathrm{F}_{b}(\widehat{\mathrm{G}})\) 的闭星子代数，它由族 \((u_{\pi})_{\pi\in\widehat{\mathrm{G}}}\) 组成，这些族满足 \(\|u_{\pi}\|\) 在无穷远处趋于 0。

令 \(\nu \in \mathcal{M}^1 (\mathrm{G})\)。对于每个 \(\pi \in \widehat{\mathrm{G}}\)，有 \(\| \pi (\nu)\| \leqslant \| \nu \|\)，所以族 \((\pi (\nu))_{\pi \in \widehat{\mathrm{G}}}\) 属于 \(\mathrm{F}_b(\widehat{\mathrm{G}})\)。

定义 1. --- 对每个测度 \(\nu \in \mathcal{M}^{1}(\mathrm{G})\)，将族 \((\pi (\nu))_{\pi \in \widehat{\mathrm{G}}}\) 视为 \(\mathrm{F}_b(\widehat{\mathrm{G}})\) 中的元素，记为 \(\overline{\mathcal{F}}_{\mathrm{G}}(\nu)\)。这样定义的从 \(\mathcal{M}^1 (\mathrm{G})\) 到 \(\mathrm{F}_b(\widehat{\mathrm{G}})\) 的映射称为 G 的傅里叶余变换，而 \(\overline{\mathcal{F}}_{\mathrm{G}}(\nu)\) 称为测度 \(\nu\) 的傅里叶余变换。

对于 \(f \in \mathrm{L}^1(\mathrm{G})\)，记 \(\overline{\mathcal{F}}_{\mathrm{G}}(f) = \overline{\mathcal{F}}_{\mathrm{G}}(f \cdot \mu)\)。

对于任意表示 \(\pi \in \widehat{\mathrm{G}}\)，映射 \(\nu \mapsto \pi(\nu)\) 是从 \(\mathcal{M}^{1}(\mathrm{G})\) 到 \(\operatorname{End}(\mathrm{E}_{\pi})\) 的含单位对合巴拿赫代数态射（V, 第 401 页引理 1）。因此，傅里叶余变换是从 \(\mathcal{M}^{1}(\mathrm{G})\) 到 \(\mathrm{F}_b(\widehat{\mathrm{G}})\) 的含单位对合巴拿赫代数态射。

